%%%%%%%%%%%%%%%%%%%%%%%%%%%%%%%%%%%%%%%%%%%%%%%%%%%%%%%%%%%%%%%%%%%%%%%%%%%%%%%%
% S6.ReKI.1 — Research Project Report
% Muhammad Sawaiz Naveed · Technische Universität Ilmenau
% Preamble follows expose.tex; additions for tables, listings and units.
%%%%%%%%%%%%%%%%%%%%%%%%%%%%%%%%%%%%%%%%%%%%%%%%%%%%%%%%%%%%%%%%%%%%%%%%%%%%%%%%

\documentclass[11pt,a4paper,titlepage=false]{scrreprt}

\usepackage[utf8]{inputenc}
\usepackage[T1]{fontenc}
\usepackage{lmodern}
\usepackage[english]{babel}
\usepackage{amsmath,amsfonts,amssymb}
\usepackage{cite}
\usepackage{hyperref}
\usepackage{graphicx}
\usepackage{url}
\usepackage{breakurl}
\usepackage{xcolor}
\usepackage{enumitem}
\usepackage[
  left=2.5cm,
  right=2.5cm,
  top=3cm,
  bottom=3cm,
]{geometry}

% --- additions beyond the expose preamble -----------------------------------
\usepackage{booktabs}      % rules for the results tables
\usepackage{longtable}     % tables spanning a page break
\usepackage{array}
\usepackage{caption}
\usepackage{float}
\usepackage{microtype}     % reduces overfull boxes in a two-column-free layout
\usepackage{listings}

\hypersetup{
    colorlinks = true,
    linkcolor  = black,
    citecolor  = black,
    urlcolor   = blue,
    pdfauthor  = {Muhammad Sawaiz Naveed},
    pdftitle   = {Agentic Verilog Testbench Generation with Pre-Simulation Error Localisation},
    pdfsubject = {S6.ReKI.1 Research Project},
    pdfkeywords= {LLM, Verilog, testbench generation, Pyverilog, error localization,
                  hardware verification, LangGraph, fault injection}
}

% Verilog listing style, used for the structural/semantic examples
\lstdefinestyle{verilog}{
  language        = Verilog,
  basicstyle      = \ttfamily\small,
  keywordstyle    = \bfseries,
  commentstyle    = \itshape\color{gray!70!black},
  numbers         = none,
  frame           = single,
  framesep        = 4pt,
  rulecolor       = \color{black!30},
  breaklines      = true,
  showstringspaces= false,
  captionpos      = b,
}
\lstset{style=verilog}

\DeclareOldFontCommand{\bf}{\normalfont\bfseries}{\mathbf}

% Tighter chapter headings: the page budget is 40 pages.
\renewcommand*{\chapterheadstartvskip}{\vspace*{-1.2\baselineskip}}

\begin{document}
\selectlanguage{english}

%%%%%%%%%%%%%%%%%%%%%%%%%%%%%%%%%%%%%%%%%%%%%%%%%%%%%%%%%%%%%%%%%%%%%%%%%%%%%%%%
%                                 Front matter
%%%%%%%%%%%%%%%%%%%%%%%%%%%%%%%%%%%%%%%%%%%%%%%%%%%%%%%%%%%%%%%%%%%%%%%%%%%%%%%%

\title{Agentic Verilog Testbench Generation\\with Pre-Simulation Error Localisation}
\author{Muhammad Sawaiz Naveed}
\date{\today}
\maketitle

\vspace{0.5cm}

\begin{tabular}{r l}
Type of work: & Research Project \\
Student:      & Muhammad Sawaiz Naveed \\
Supervisor:   & Bing Wen \\
Topic ID:     & S6.ReKI.1 \\
Department:   & Technische Universität Ilmenau \\
Semester:     & Summer Semester 2025/2026 \\
Deadline:     & September 1, 2026 \\
\end{tabular}

\vspace{4mm}

\begin{tabular}{l p{0.75\textwidth}}
\textbf{Keywords:} & Large Language Models; Verilog; Testbench Generation; Pyverilog;
Error Localisation; Hardware Verification; LangGraph; Fault Injection \\
\end{tabular}

\vspace{8mm}

\begin{abstract}
% PHASE 6 — written last, after the conclusion is fixed.
\textit{[Abstract placeholder. Written in Phase 6.]}
\end{abstract}

\vspace{4mm}
\tableofcontents

%%%%%%%%%%%%%%%%%%%%%%%%%%%%%%%%%%%%%%%%%%%%%%%%%%%%%%%%%%%%%%%%%%%%%%%%%%%%%%%%
\chapter{Introduction}\label{ch:intro}
% PHASE 6 · target 3.0 pp
\section{Verification and the Cost of Testbenches}\label{sec:intro-cost}
\section{How LLM-Generated Testbenches Fail}\label{sec:intro-failure}
\section{Hypothesis and Research Questions}\label{sec:intro-rq}
\section{Summary of Findings}\label{sec:intro-findings}
\section{Contributions}\label{sec:intro-contrib}
\section{Structure of this Report}\label{sec:intro-roadmap}

%%%%%%%%%%%%%%%%%%%%%%%%%%%%%%%%%%%%%%%%%%%%%%%%%%%%%%%%%%%%%%%%%%%%%%%%%%%%%%%%
\chapter{Background and Related Work}\label{ch:background}
% PHASE 5 · target 5.5 pp
\section{Simulation-Based Verification}\label{sec:bg-verification}
\section{Structural and Semantic Testbench Faults}\label{sec:bg-fault-classes}
\section{Large Language Models for HDL}\label{sec:bg-llm-hdl}
\section{Static Analysis of Verilog}\label{sec:bg-static}
\section{Evaluation Criteria: Eval0, Eval1 and Eval2}\label{sec:bg-eval}
\section{AutoBench}\label{sec:bg-autobench}
\section{Further Related Work}\label{sec:bg-related}
\section{Gap Addressed by this Work}\label{sec:bg-gap}

%%%%%%%%%%%%%%%%%%%%%%%%%%%%%%%%%%%%%%%%%%%%%%%%%%%%%%%%%%%%%%%%%%%%%%%%%%%%%%%%
\chapter{System Design}\label{ch:design}
% PHASE 1 · target 5.5 pp
\section{Why a Graph Rather than a Script}\label{sec:design-graph}
\section{Pipeline Topology and State}\label{sec:design-topology}
\section{Generation Stages}\label{sec:design-generation}
\section{Deterministic Standardisation}\label{sec:design-standardiser}
\section{The Repair Loop}\label{sec:design-repair}
\section{Evaluation Harness}\label{sec:design-harness}
\section{Engineering Practices}\label{sec:design-practices}

%%%%%%%%%%%%%%%%%%%%%%%%%%%%%%%%%%%%%%%%%%%%%%%%%%%%%%%%%%%%%%%%%%%%%%%%%%%%%%%%
\chapter{The Static Localiser}\label{ch:localiser}
% PHASE 1 · target 4.0 pp
\section{Design Principle and Scope}\label{sec:loc-principle}
\section{The Six Checks}\label{sec:loc-checks}
\section{Width Mismatch: Silent at Simulation Time}\label{sec:loc-width}
\section{A Check that was Removed}\label{sec:loc-removed}
\section{False Positives as the Primary Risk}\label{sec:loc-falsepos}
\section{Parse Robustness on Generated Code}\label{sec:loc-parse}

%%%%%%%%%%%%%%%%%%%%%%%%%%%%%%%%%%%%%%%%%%%%%%%%%%%%%%%%%%%%%%%%%%%%%%%%%%%%%%%%
\chapter{Validating the Localiser}\label{ch:validation}
% PHASE 2 · target 4.0 pp
\section{Why Pass Rates Cannot Answer RQ2}\label{sec:val-motivation}
\section{Fault Injection Method}\label{sec:val-method}
\section{Injectors and Negative Controls}\label{sec:val-injectors}
\section{Methodological Safeguards}\label{sec:val-safeguards}
\section{Detection, Localisation and False Positives}\label{sec:val-results}
\section{The Case Only Static Analysis Detects}\label{sec:val-unobserved}
\section{Boundaries of the Method}\label{sec:val-limits}

%%%%%%%%%%%%%%%%%%%%%%%%%%%%%%%%%%%%%%%%%%%%%%%%%%%%%%%%%%%%%%%%%%%%%%%%%%%%%%%%
\chapter{Experimental Design}\label{ch:experiment}
% PHASE 2 · target 3.0 pp
\section{Three Sweeps}\label{sec:exp-sweeps}
\section{Circuit Sets and Selection Criteria}\label{sec:exp-circuits}
\section{Ablation Arms}\label{sec:exp-arms}
\section{The No-Diagnostics Control}\label{sec:exp-control}
\section{Prompt Freezing}\label{sec:exp-freeze}
\section{Statistical Treatment}\label{sec:exp-stats}

%%%%%%%%%%%%%%%%%%%%%%%%%%%%%%%%%%%%%%%%%%%%%%%%%%%%%%%%%%%%%%%%%%%%%%%%%%%%%%%%
\chapter{Results}\label{ch:results}
% PHASE 3 · target 6.5 pp
\section{Pipeline Reliability}\label{sec:res-reliability}
\section{RQ1: Where Failures Originate}\label{sec:res-rq1}
\section{RQ2a: Detection Capability}\label{sec:res-rq2a}
\section{RQ2b: Yield on Generated Testbenches}\label{sec:res-rq2b}
\section{RQ3: Ablation Across Feedback Strategies}\label{sec:res-rq3}
\section{The Measured Variance Floor}\label{sec:res-variance}
\section{Uninformed Retry Degrades Compilation}\label{sec:res-retry}
\section{RQ4: Cost Against Quality}\label{sec:res-rq4}
\section{Consistency Across Models and Circuit Sets}\label{sec:res-consistency}

%%%%%%%%%%%%%%%%%%%%%%%%%%%%%%%%%%%%%%%%%%%%%%%%%%%%%%%%%%%%%%%%%%%%%%%%%%%%%%%%
\chapter{Discussion}\label{ch:discussion}
% PHASE 4 · target 5.5 pp
\section{Excluding the Alternative Explanations}\label{sec:disc-alternatives}
\section{Why Generated Testbenches Are Structurally Sound}\label{sec:disc-why}
\section{The Same Technique in a Different Model Generation}\label{sec:disc-autobench}
\section{Where Static Analysis Retains Value}\label{sec:disc-value}
\section{On Measurement Practice}\label{sec:disc-measurement}
\section{Threats to Validity}\label{sec:disc-threats}
\section{Corrections Made During the Work}\label{sec:disc-corrections}

%%%%%%%%%%%%%%%%%%%%%%%%%%%%%%%%%%%%%%%%%%%%%%%%%%%%%%%%%%%%%%%%%%%%%%%%%%%%%%%%
\chapter{Conclusion and Future Work}\label{ch:conclusion}
% PHASE 6 · target 2.0 pp
\section{Answers to the Research Questions}\label{sec:concl-rq}
\section{Contributions in Retrospect}\label{sec:concl-contrib}
\section{Future Work}\label{sec:concl-future}

%%%%%%%%%%%%%%%%%%%%%%%%%%%%%%%%%%%%%%%%%%%%%%%%%%%%%%%%%%%%%%%%%%%%%%%%%%%%%%%%
%                                 References
%%%%%%%%%%%%%%%%%%%%%%%%%%%%%%%%%%%%%%%%%%%%%%%%%%%%%%%%%%%%%%%%%%%%%%%%%%%%%%%%

{\small
\begingroup
\let\chapter\section
\bibliographystyle{ieeetr}

\begin{thebibliography}{99}

\bibitem{wilson2022}
H.~Foster,
``The 2022 Wilson Research Group Functional Verification Study,''
Siemens, Tech. Rep., Dec. 2022.

\bibitem{liu2023verilogeval}
M.~Liu, N.~Pinckney, B.~Khailany, and H.~Ren,
``VerilogEval: Evaluating Large Language Models for Verilog Code Generation,''
in \textit{Proc. IEEE/ACM Int. Conf. Computer-Aided Design (ICCAD)}, 2023.

\bibitem{qiu2024autobench}
R.~Qiu, G.~L.~Zhang, R.~Drechsler, U.~Schlichtmann, and B.~Li,
``AutoBench: Automatic Testbench Generation and Evaluation Using LLMs for HDL Design,''
in \textit{Proc. ACM/IEEE Int. Symp. Machine Learning for CAD (MLCAD)}, 2024.
arXiv:2407.03891.

\bibitem{takamaeda2015pyverilog}
S.~Takamaeda-Yamazaki,
``Pyverilog: A Python-Based Hardware Design Processing Toolkit for Verilog HDL,''
in \textit{Applied Reconfigurable Computing (ARC)}, 2015.

\bibitem{thakur2023autochip}
S.~Thakur, J.~Blocklove, H.~Pearce, B.~Tan, S.~Garg, and R.~Karri,
``AutoChip: Automating HDL Generation Using LLM Feedback,''
arXiv:2311.04887, 2023.

\bibitem{blocklove2023chipchat}
J.~Blocklove, S.~Garg, R.~Karri, and H.~Pearce,
``Chip-Chat: Challenges and Opportunities in Conversational Hardware Design,''
in \textit{Proc. ACM/IEEE Int. Symp. Machine Learning for CAD (MLCAD)}, 2023.

\bibitem{chang2023chipgpt}
K.~Chang, Y.~Wang, H.~Ren, M.~Wang, S.~Liang, Y.~Han, H.~Li, and X.~Li,
``ChipGPT: How Far Are We From Natural Language Hardware Design,''
arXiv:2305.14019, 2023.

\bibitem{orenes2023llmformal}
M.~Orenes-Vera, M.~Martonosi, and D.~Wentzlaff,
``Using LLMs to Facilitate Formal Verification of RTL,''
arXiv:2309.09437, 2023.

\bibitem{zhang2023llm4dv}
Z.~Zhang, G.~Chadwick, H.~McNally, Y.~Zhao, and R.~Mullins,
``LLM4DV: Using Large Language Models for Hardware Test Stimuli Generation,''
arXiv:2310.04535, 2023.

\bibitem{wilson1927}
E.~B.~Wilson,
``Probable Inference, the Law of Succession, and Statistical Inference,''
\textit{Journal of the American Statistical Association}, vol.~22, no.~158,
pp.~209--212, 1927.

\bibitem{fisher1922}
R.~A.~Fisher,
``On the Interpretation of $\chi^2$ from Contingency Tables, and the Calculation of P,''
\textit{Journal of the Royal Statistical Society}, vol.~85, no.~1, pp.~87--94, 1922.

\bibitem{langgraph}
LangChain Inc.,
``LangGraph: Building Stateful, Multi-Actor Applications with LLMs,'' 2024.
\url{https://langchain-ai.github.io/langgraph/}

\bibitem{iverilog}
S.~Williams,
``Icarus Verilog,'' open-source Verilog simulator.
\url{https://steveicarus.github.io/iverilog/}

\bibitem{verible}
Google,
``Verible: SystemVerilog Parser, Style-Linter and Formatter.''
\url{https://github.com/chipsalliance/verible}

\bibitem{jia2011mutation}
Y.~Jia and M.~Harman,
``An Analysis and Survey of the Development of Mutation Testing,''
\textit{IEEE Trans. Software Engineering}, vol.~37, no.~5, pp.~649--678, 2011.

\end{thebibliography}
\endgroup
}

\end{document}
